% Report for the NLP course project P6 The Apprentice Model
% Based on the Springer Nature article template

\documentclass[pdflatex,sn-mathphys-num]{sn-jnl}

\usepackage{graphicx}
\usepackage{amsmath,amssymb,amsfonts}
\usepackage{booktabs}
\usepackage{xcolor}
\usepackage{textcomp}
\usepackage{url}
\usepackage{listings}


\raggedbottom

\begin{document}

\title[The Apprentice Model for Toxicity Detection]{Distilling an LLM Teacher into a Compact Transformer for Toxic Comment Classification}

\author{\fnm{Vadim} \sur{Sokolov}}

\affil{\orgdiv{Master's Degree in Computer Science}, \orgname{Universit\`a degli Studi di Milano}, \orgaddress{\city{Milan}, \country{Italy}}}

\abstract{This project studies knowledge distillation for Natural Language Processing through an ``apprentice model''. 
The goal is not to investigate if a compact student model can imitate the decision behaviour of a larger teacher model on a specific text classification task.
The selected task is toxic comment classification. 
The experimental pipeline compares a classical TF-IDF and Logistic Regression baseline, a supervised BERT tiny student model, that was trained on original dataset labels, and an LLM OpenAI teacher used to relabel a controlled subset of examples. 
The results show that the supervised BERT-tiny student improves over the TF-IDF baseline in accuracy and F1 score, but is larger and slower at inference time. 
The teacher-label analysis also shows meaningful disagreements between dataset labels and the LLM teacher. It shows that the teacher is utilizing a stricter decision boundary for hostile or demeaning language.
}

\keywords{Natural Language Processing, Knowledge Distillation, Toxicity Detection, BERT, TF-IDF, Large Language Models, GPT}

\maketitle

\section{Introduction}\label{sec:introduction}


\end{document}
